% ============================================================
%  Geometría Analítica · LM · Universidad de Sonora
%  Semana 3 — Ángulos, distancias, familias de rectas y
%             puntos notables del triángulo
%  Semana SIN laboratorio (el laboratorio es quincenal;
%  la Construcción 2 va la próxima semana)
%  Compilar en Overleaf con pdfLaTeX
% ============================================================
\documentclass[aspectratio=169,10pt]{beamer}
\usepackage[utf8]{inputenc}
\usepackage[T1]{fontenc}
\usepackage[spanish,es-nodecimaldot,es-noquoting]{babel}
\usepackage{lmodern}
\usepackage{amsmath,amssymb}
\usepackage{booktabs}
\usepackage{tikz}
\usetikzlibrary{arrows.meta,calc}
\usepackage{graphicx}
\graphicspath{{figuras/}}
% ---------- Tema y colores ----------
\usetheme{default}
\setbeamertemplate{navigation symbols}{}
\definecolor{turquesaGA}{HTML}{0E7C8C}
\definecolor{turquesaClaro}{HTML}{E2F2F4}
\definecolor{grisT}{HTML}{5F5E5A}
\definecolor{negroP}{HTML}{2B2B2B}
\definecolor{terracota}{HTML}{C1502E}
\definecolor{dorado}{HTML}{B07C10}
\definecolor{verdeOk}{HTML}{1D9E75}
\definecolor{rojoAc}{HTML}{C0392B}
\setbeamercolor{structure}{fg=turquesaGA}
\setbeamercolor{frametitle}{fg=white,bg=turquesaGA}
\setbeamercolor{title}{fg=turquesaGA}
\setbeamerfont{frametitle}{size=\large,series=\bfseries}
\setbeamertemplate{frametitle}{%
  \nointerlineskip
  \begin{beamercolorbox}[wd=\paperwidth,ht=2.6ex,dp=1.2ex,leftskip=0.3cm]{frametitle}
    \usebeamerfont{frametitle}\insertframetitle
  \end{beamercolorbox}}
\setbeamertemplate{footline}{%
  \hfill\textcolor{grisT}{\tiny Geometría Analítica · LM · Unison \quad
  \insertframenumber/\inserttotalframenumber}\hspace{0.3cm}\vspace{0.15cm}}
\setbeamertemplate{itemize item}{\textcolor{turquesaGA}{\raisebox{0.15ex}{\tiny$\blacksquare$}}}
\setbeamertemplate{itemize subitem}{\textcolor{grisT}{\raisebox{0.15ex}{\tiny$\bullet$}}}
% ---------- Caja auxiliar para destacar ideas ----------
\newcommand{\destaca}[1]{%
  \par\vspace{3pt}\noindent
  \colorbox{turquesaClaro}{%
    \parbox{\dimexpr\linewidth-2\fboxsep\relax}{\vspace{3pt}#1\vspace{3pt}}}%
  \par\vspace{3pt}}
% ---------- Hojas de separación por sesión ----------
%  Las sesiones no llevan número: se identifican por su título.
\newcommand{\hojaSesion}[2]{%
  \vfill
  \begin{beamercolorbox}[wd=\textwidth,sep=16pt,center]{frametitle}
    {\Large\bfseries #1}\\[11pt]
    \begin{minipage}{0.82\textwidth}\centering\small #2\end{minipage}
  \end{beamercolorbox}
  \vfill}
% ---------- Abreviaturas ----------
\newcommand{\R}{\mathbb{R}}
\newcommand{\pp}[2]{\langle #1,#2\rangle}
\newcommand{\norma}[1]{\lVert #1\rVert}
% ============================================================
\title{\LARGE\textbf{Geometría Analítica}}
\subtitle{Semana 3 · Ángulos, distancias y puntos notables del triángulo}
\author{Prof. Rosalía Gpe. Hernández Amador}
\institute{Departamento de Matemáticas · Universidad de Sonora}
\date{Semestre 2026-2}
\begin{document}
% ============================================================
\begin{frame}[plain]
  \titlepage
\end{frame}
% ============================================================
\begin{frame}{Ruta de la semana}
  \centering
  {\small
  \begin{tabular}{@{}l@{}}
    \toprule
    \textbf{Las cinco sesiones de la semana, en orden} \\
    \midrule
    Ángulo entre dos rectas, por dos métodos \\
    Proyección de un vector; distancia de un punto a una recta \\
    Taller: la forma normal clásica; tarea de redacción \\
    Familias de rectas con un parámetro \\
    Taller: los puntos notables del triángulo \\
    \bottomrule
  \end{tabular}}

  \vspace{0.4cm}
  \raggedright
  \destaca{Esta semana el producto punto de la semana pasada medirá el
  ángulo entre rectas y producirá \textbf{todas} las distancias del
  curso. Al final juntamos las herramientas de estas tres semanas en
  un problema clásico, los puntos notables del triángulo.
  \textbf{Esta semana no hay laboratorio} (es quincenal); la
  Construcción 2 va la próxima semana.}
\end{frame}
% ============================================================
\section{Ángulo entre dos rectas}
% ---------- hoja de sesión ----------
\begin{frame}[plain]
  \hojaSesion{Ángulo entre dos rectas}{%
  Dos métodos para el mismo ángulo: el vectorial, con producto punto
  de directores, y el clásico, con pendientes. Comparación de sus
  alcances (las rectas verticales distinguen a uno del otro) y un
  ejemplo resuelto por ambas vías, con la coincidencia como
  verificador.}
\end{frame}
\begin{frame}{Pregunta pendiente}
  Cerramos la semana 2 con esta pregunta. La ecuación $5x-2y=20$
  contiene, a la vista, un vector perpendicular a la recta y sus dos
  cortes con los ejes. ¿Cuáles son?

  \vspace{0.3cm}
  \begin{itemize}
    \item El vector normal se lee de los coeficientes: $n=(5,-2)$.
    \item Con $y=0$ resulta $x=4$; con $x=0$ resulta $y=-10$.
      Los cortes son $(4,0)$ y $(0,-10)$.
  \end{itemize}

  \vspace{0.3cm}
  \destaca{Dividiendo entre $20$ obtenemos la forma simétrica
  $\dfrac{x}{4}+\dfrac{y}{-10}=1$, donde los cortes quedan a la vista.
  Nótese la moraleja: \textbf{cada forma de la ecuación exhibe un dato
  distinto}, y reescribir la ecuación es una manera de leer datos
  nuevos.}
\end{frame}
% ============================================================
\begin{frame}{Método 1: con vectores directores}
  Si $d^{(1)}$ y $d^{(2)}$ son vectores directores de las rectas,
  \[
    \cos\theta=\frac{\bigl|\pp{d^{(1)}}{d^{(2)}}\bigr|}
      {\norma{d^{(1)}}\,\norma{d^{(2)}}}
  \]
  (el valor absoluto selecciona el ángulo \textit{agudo} entre las
  rectas).

  \vspace{0.4cm}
  \destaca{\textbf{Ventajas de este método:} es válido sin
  excepciones, incluidas las rectas verticales; y se extiende sin
  cambios a $\R^3$, donde las pendientes ni siquiera existen. Cuando
  lleguemos al espacio, en la Unidad 3, esta será la única fórmula
  que necesitaremos recordar.}
\end{frame}
% ============================================================
\begin{frame}{Método 2: con pendientes (fórmula clásica)}
  Si las pendientes son $m_1$ y $m_2$ y las rectas no son
  perpendiculares,
  \[
    \tan\theta=\left|\frac{m_2-m_1}{1+m_1m_2}\right| .
  \]
  Resulta cómoda cuando las pendientes ya se conocen.

  \vspace{0.3cm}
  \begin{center}
  {\small
  \begin{tabular}{@{}lcc@{}}
    \toprule
    & \textbf{con pendientes} & \textbf{con vectores} \\
    \midrule
    paralelas & $m_1=m_2$ & $d^{(1)}=\lambda\, d^{(2)}$ \\
    perpendiculares & $m_1\,m_2=-1$ & $\pp{d^{(1)}}{d^{(2)}}=0$ \\
    \bottomrule
  \end{tabular}}
  \end{center}

  \vspace{0.3cm}
  \destaca{\textbf{Observación importante.} La condición $m_1m_2=-1$
  no es aplicable cuando una recta es vertical y la otra horizontal,
  aunque son perpendiculares. La condición con producto punto no tiene
  excepciones. Tener dos lenguajes y saber cuál falla dónde también
  es contenido del curso.}
\end{frame}
% ============================================================
\begin{frame}{Ejemplos}
  \textbf{Ejemplo 1.} Determinar el ángulo entre las rectas $y=2x-1$
  y $y=-\tfrac13 x+4$, por ambos métodos.

  \vspace{0.1cm}

  {\small \textit{Solución.} Con pendientes:
  $\tan\theta=\bigl|\tfrac{-1/3-2}{1+2\cdot(-1/3)}\bigr|
  =\bigl|\tfrac{-7/3}{1/3}\bigr|=7$, de donde
  $\theta\approx 81.9^\circ$. Con directores $(1,2)$ y $(3,-1)$:
  $\cos\theta=\tfrac{|3-2|}{\sqrt5\,\sqrt{10}}=\tfrac{1}{\sqrt{50}}$,
  que produce el mismo $\theta\approx 81.9^\circ$.}

  \vspace{0.3cm}

  \textbf{Ejemplo 2.} ¿Son perpendiculares $3x+2y=7$ y
  $X=(1,1)+t\,(2,-3)$?

  \vspace{0.1cm}

  {\small \textit{Solución.} El normal de la primera es $(3,2)$; su
  director es, por tanto, $(-2,3)$. Como
  $\pp{(-2,3)}{(2,-3)}=-13\neq 0$, no son perpendiculares. De hecho
  $(2,-3)=-(-2,3)$: las rectas son \textit{paralelas}.}

  \vspace{0.2cm}
  \destaca{\small Nótese, en el Ejemplo 1, que ambos métodos
  coinciden. Resolver por dos vías es el verificador de cálculo más
  eficaz que tenemos, y así lo pediremos en los talleres.}
\end{frame}
% ============================================================
\section{Proyección y distancia punto-recta}
% ---------- hoja de sesión ----------
\begin{frame}[plain]
  \hojaSesion{Proyección y distancia punto-recta}{%
  La proyección, es decir, la componente de un vector en la dirección
  de otro. Con ella, el teorema de la distancia de un punto a una
  recta, primero con normal unitario (donde la idea es transparente)
  y después en el caso general, que produce la fórmula clásica.}
\end{frame}
\begin{frame}{Proyección de un vector sobre otro}
  \begin{columns}[T]
    \begin{column}{0.52\textwidth}
      Pregunta: ¿qué tanto de $u$ apunta en la dirección de $v$?

      \vspace{0.2cm}
      \destaca{\textbf{Proyección.} La \textbf{proyección escalar} de
      $u$ sobre $v$ es
      \[
        \frac{\pp{u}{v}}{\norma{v}}\;=\;\norma{u}\cos\theta ,
      \]
      y el \textbf{vector proyección} es ese escalar multiplicado por
      el vector unitario en la dirección de $v$.}

      \vspace{0.2cm}
      Si $\norma{v}=1$, la proyección escalar es simplemente
      $\pp{u}{v}$. Con vectores unitarios,
      \textbf{proyectar se reduce a multiplicar}.
    \end{column}
    \begin{column}{0.44\textwidth}
      \begin{center}
      \begin{tikzpicture}[scale=.75,>=Stealth]
        \draw[very thin,color=gray!30] (-0.9,-0.9) grid (4.1,2.9);
        \draw[->] (-0.9,0) -- (4.1,0);
        \draw[->] (0,-0.9) -- (0,2.9);
        \draw[dorado,very thick,->] (0,0) -- (3.6,0.9)
          node[right]{\small $v$};
        \draw[terracota,very thick,->] (0,0) -- (1.4,2.3)
          node[above]{\small $u$};
        \coordinate (S) at ({(1.4*3.6+2.3*0.9)/(3.6*3.6+0.9*0.9)*3.6},
                            {(1.4*3.6+2.3*0.9)/(3.6*3.6+0.9*0.9)*0.9});
        \draw[dashed,gray,thick] (1.4,2.3) -- (S);
        \draw[turquesaGA,ultra thick] (0,0) -- (S)
          node[midway,below,sloped]{\small proyección};
      \end{tikzpicture}
      \end{center}
    \end{column}
  \end{columns}
\end{frame}
% ============================================================
\begin{frame}{El teorema de la distancia punto-recta}
  \destaca{\textbf{Teorema.} Si la recta $\ell$ tiene ecuación
  $Ax+By=C$ con vector normal \textbf{unitario} $n=(A,B)$ (es decir,
  $A^2+B^2=1$), entonces para cualquier punto $Q=(q_1,q_2)$:
  \[
    d(Q,\ell)=\bigl|\,\pp{n}{Q}-C\,\bigr|
    =\bigl|\,Aq_1+Bq_2-C\,\bigr| .
  \]}

  \vspace{0.1cm}
  \textit{Idea de la demostración:} $\pp{n}{Q}$ es la proyección de
  $Q$ en la dirección perpendicular a la recta (su ``nivel'') y $C$
  es el nivel común de los puntos de la recta. La distancia es la
  diferencia de niveles.

  \vspace{0.3cm}
  \destaca{\textbf{Caso general.} Si el normal no es unitario,
  dividimos entre su norma:
  \[
    d(Q,\ell)=\frac{|Aq_1+Bq_2-C|}{\sqrt{A^2+B^2}} .
  \]
  Esta es la fórmula clásica de los libros de texto; ahora sabemos de
  dónde sale.}
\end{frame}
% ============================================================
\begin{frame}{Ejemplos}
  \textbf{Ejemplo 3.} Distancia del punto $(2,7)$ a la recta
  $3x+4y=10$.

  \vspace{0.1cm}

  {\small \textit{Solución.}
  $d=\dfrac{|3\cdot 2+4\cdot 7-10|}{\sqrt{9+16}}=\dfrac{24}{5}$.}

  \vspace{0.3cm}

  \textbf{Ejemplo 4.} Distancia del origen a la misma recta.

  \vspace{0.1cm}

  {\small \textit{Solución.} $d=\dfrac{|0+0-10|}{5}=2$. Para el
  origen, la fórmula se reduce a $|C|/\norma{n}$.}

  \vspace{0.3cm}

  \textbf{Pregunta.} ¿Para qué puntos $Q$ la fórmula produce cero?

  \vspace{0.2cm}
  \destaca{\small \textit{Respuesta.} Exactamente para los puntos de
  la recta: distancia cero equivale a satisfacer la ecuación. Son las
  dos caras del lugar geométrico (\textit{todos} y
  \textit{solamente}) una vez más.}
\end{frame}
% ============================================================
\section{Taller: la forma normal}
% ---------- hoja de sesión ----------
\begin{frame}[plain]
  \hojaSesion{Taller: la forma normal}{%
  La forma normal de la recta en la notación clásica de Lehmann y su
  traducción al lenguaje vectorial: son la misma fórmula. El mismo
  problema de distancia resuelto por las dos vías, y la tarea de
  redacción de la semana.}
\end{frame}
\begin{frame}{La forma normal de la recta}
  En el texto de Lehmann, la recta aparece escrita en la
  \textit{forma normal}:
  \[
    x\cos\omega+y\,\mathrm{sen}\,\omega-p=0 ,
  \]
  donde $\omega$ es el ángulo del vector normal y $p$ la distancia de
  la recta al origen.

  \vspace{0.3cm}

  Traducción al lenguaje de esta semana:
  \[
    (\cos\omega,\ \mathrm{sen}\,\omega) = \text{el normal unitario},
    \qquad p = \text{el nivel de la recta.}
  \]

  \vspace{0.1cm}
  \destaca{No se trata de una fórmula nueva. Es la misma expresión
  $\pp{n}{X}=p$ con $\norma{n}=1$, en la notación clásica. Saber
  \textbf{traducir entre notaciones} es parte de la formación
  matemática: no todos los libros usan la misma notación, pero dicen
  lo mismo.}
\end{frame}
% ============================================================
\begin{frame}{Taller: el mismo problema por dos vías}
  Distancia de $(1,5)$ a la recta que pasa por $(0,2)$ y $(4,0)$:

  \vspace{0.3cm}
  \begin{columns}[T]
    \begin{column}{0.48\textwidth}
      \textbf{Vía clásica:} pendiente, forma punto--pendiente, forma
      general, división entre $\sqrt{A^2+B^2}$, sustitución.
    \end{column}
    \begin{column}{0.48\textwidth}
      \textbf{Vía vectorial:} director $(4,-2)$; normal $(2,4)$,
      simplificado a $(1,2)$; nivel $C=\pp{(1,2)}{(0,2)}=4$;
      \[
        d=\frac{|1\cdot 1+2\cdot 5-4|}{\sqrt{5}}=\frac{7}{\sqrt 5}.
      \]
    \end{column}
  \end{columns}

  \vspace{0.4cm}
  \destaca{Ambas vías deben coincidir. Si no coinciden, hay un error
  que localizar, y aprender a \textbf{localizar errores propios} es
  una habilidad central de la formación matemática. Nótese que la raíz
  la dejamos indicada: $7/\sqrt{5}$ es la respuesta exacta.}
\end{frame}
% ============================================================
\begin{frame}{Tarea de redacción}
  \destaca{\textbf{Tarea de redacción (se entrega en Moodle).}
  Explicar, con un argumento completo y una figura, por qué
  $d(Q,\ell)=|\pp{n}{Q}-C|$ cuando $\norma{n}=1$. El texto debe
  indicar: (a) qué representa $\pp{n}{Q}$; (b) por qué el punto de
  $\ell$ más cercano a $Q$ es el pie de la perpendicular; (c) en qué
  paso se utiliza que $\norma{n}=1$.}

  \vspace{0.3cm}
  En matemáticas no basta con llegar al resultado; hay que saber
  \textbf{comunicarlo} de modo que otra persona pueda seguir el
  argumento sin ayuda de quien lo escribió. Esta tarea evalúa
  exactamente eso.

  \vspace{0.3cm}
  \destaca{\small \textbf{Aviso de laboratorio.} Esta semana no hay
  laboratorio (es quincenal). La \textbf{Construcción 2} (el pie de
  la perpendicular, construido y no dibujado) la trabajaremos la
  próxima semana, y esta tarea de redacción es su mejor preparación.}
\end{frame}
% ============================================================
\section{Familias de rectas}
% ---------- hoja de sesión ----------
\begin{frame}[plain]
  \hojaSesion{Familias de rectas con un parámetro}{%
  Una ecuación con un parámetro no describe una recta sino una familia
  completa: el haz de rectas por un punto, las familias de paralelas,
  y la familia que pasa por la intersección de dos rectas dadas (una
  técnica que resuelve problemas sin calcular la intersección).}
\end{frame}
\begin{frame}{Familias de rectas}
  \begin{columns}[T]
    \begin{column}{0.52\textwidth}
      \textbf{Haz de rectas por un punto.} Al variar $m$ en
      \[
        y-3=m\,(x-2)
      \]
      se obtienen \textbf{todas} las rectas que pasan por $(2,3)$
      \ldots{} salvo una. ¿Cuál falta?

      \vspace{0.2cm}

      \textbf{Familia de paralelas.} Al variar $C$ en
      \[
        x+2y=C
      \]
      se obtienen todas las rectas con normal $(1,2)$: una familia de
      paralelas que barre el plano.
    \end{column}
    \begin{column}{0.44\textwidth}
      \begin{center}
      \begin{tikzpicture}[scale=.55,>=Stealth]
        \draw[very thin,color=gray!30] (-0.9,-0.9) grid (4.9,4.1);
        \draw[->] (-0.9,0) -- (5.1,0);
        \draw[->] (0,-0.9) -- (0,4.3);
        \draw[terracota,thick] (-0.5,0.5) -- (3.1,4.1);
        \draw[terracota!70,thick] (-0.9,1.55) -- (4.2,4.1);
        \draw[terracota!45,thick] (-0.9,3.725) -- (4.9,2.275);
        \draw[gray!70,thick,dashed] (2,-0.9) -- (2,4.1);
        \fill[turquesaGA] (2,3) circle(3pt);
      \end{tikzpicture}\\[2pt]
      {\scriptsize El haz por $(2,3)$; la vertical, punteada,
      no la produce ningún $m$.}
      \end{center}
    \end{column}
  \end{columns}

  \vspace{0.1cm}
  \destaca{La que falta en el haz es la \textbf{vertical} $x=2$: no
  tiene pendiente, así que ningún valor de $m$ la produce. El
  parámetro funciona como un \textbf{selector}. Cada valor elige a un
  miembro de la familia, y elegir bien el valor es resolver el
  problema.}
\end{frame}
% ============================================================
\begin{frame}{Seleccionar un miembro de la familia}
  \textbf{Ejemplo 5.} ¿Qué recta de la familia $x+2y=C$ pasa por el
  punto $(4,1)$?

  \vspace{0.1cm}

  {\small \textit{Solución.} Imponemos la condición: $C=4+2\cdot 1=6$.
  La recta es $x+2y=6$.}

  \vspace{0.3cm}

  \textbf{Ejemplo 6.} ¿Qué recta del haz $y-3=m\,(x-2)$ pasa por el
  origen?

  \vspace{0.1cm}

  {\small \textit{Solución.} Con $(0,0)$: $-3=m\,(-2)$, es decir
  $m=\tfrac32$. La recta es $y=\tfrac32 x$. Verificación: pasa por
  $(2,3)$, pues $\tfrac32\cdot 2=3$, y por el origen. \checkmark}

  \vspace{0.3cm}
  \destaca{El esquema es siempre el mismo: la familia aporta la
  \textbf{forma general} de la ecuación y la condición extra
  \textbf{despeja el parámetro}. Compárese con el método de los
  lugares geométricos; es la misma idea, con un solo grado de
  libertad.}
\end{frame}
% ============================================================
\begin{frame}{La familia por la intersección de dos rectas}
  Dadas $\ell_1:\,3x+y-5=0$ y $\ell_2:\,x-y+1=0$, considérese la
  familia
  \[
    (3x+y-5)+k\,(x-y+1)=0 .
  \]
  Todo punto que satisface \textbf{ambas} ecuaciones satisface la
  combinación, para cualquier $k$, así que la familia entera pasa por
  el punto de intersección de $\ell_1$ y $\ell_2$.

  \vspace{0.2cm}

  \textbf{Ejemplo 7.} Determinar la recta que pasa por la intersección
  de $\ell_1$ y $\ell_2$ y por el origen, \textit{sin calcular la
  intersección}.

  \vspace{0.1cm}

  {\small \textit{Solución.} Con $(0,0)$: $-5+k\,(1)=0$, es decir
  $k=5$. La combinación queda $8x-4y=0$, o sea $y=2x$.
  Verificación por la otra vía: la intersección es $(1,2)$ (resolviendo
  el sistema), y en efecto $y=2x$ pasa por $(1,2)$ y por el origen.
  \checkmark}

  \vspace{0.2cm}
  \destaca{\small La técnica evita calcular la intersección, y la
  verificación la calculó de todos modos: \textbf{dos vías, un
  resultado}. Obsérvese que ningún valor de $k$ produce a $\ell_2$
  misma; la familia la pierde, igual que el haz pierde a la vertical.
  En la sección \textit{El invernadero} del libro hay un laboratorio
  interactivo para mover $k$.}
\end{frame}
% ============================================================
\section{Taller: puntos notables del triángulo}
% ---------- hoja de sesión ----------
\begin{frame}[plain]
  \hojaSesion{Taller: los puntos notables del triángulo}{%
  Las herramientas de tres semanas (mediatrices, punto medio con
  vectores, perpendicularidad con producto punto) se juntan en un
  problema clásico: circuncentro, baricentro y ortocentro de un
  triángulo concreto, y la sorpresa de que los tres quedan alineados.}
\end{frame}
\begin{frame}{Tres centros para un mismo triángulo}
  \centering
  {\small
  \begin{tabular}{@{}lll@{}}
    \toprule
    \textbf{Punto} & \textbf{Donde se cruzan} & \textbf{Propiedad} \\
    \midrule
    Circuncentro & las tres mediatrices & equidista de los \textbf{vértices} \\
    Baricentro & las tres medianas & centro de gravedad; divide $2{:}1$ \\
    Ortocentro & las tres alturas & pie perpendicular desde cada vértice \\
    \bottomrule
  \end{tabular}}
  \par\vspace{0.4cm}
  \raggedright
  Que las tres rectas de cada tipo se corten en \textbf{un solo} punto
  no es evidente (son teoremas de la geometría euclidiana). Hoy no los
  demostraremos en general; los \textbf{verificaremos} en un triángulo
  concreto, con coordenadas, y de paso practicaremos todo lo visto en
  estas tres semanas.

  \vspace{0.2cm}
  \destaca{Cada punto notable usa una herramienta de la semana 2: la
  \textbf{mediatriz}, el \textbf{punto medio} escrito con vectores y
  la \textbf{perpendicularidad} con producto punto.}
\end{frame}
% ============================================================
\begin{frame}{Taller: el triángulo $A=(0,0)$, $B=(6,0)$, $C=(2,6)$}
  Por equipos, calculamos:

  \vspace{0.2cm}
  \begin{enumerate}
    \item El \textbf{circuncentro}: intersección de dos mediatrices
      (la tercera es la verificación).
    \item El \textbf{baricentro}: con $M$ el punto medio de $BC$,
      \[
        G = A+\tfrac{2}{3}\,(M-A)
      \]
      (nótese la escritura: punto más vector, como el punto medio de la semana 2).
    \item El \textbf{ortocentro}: intersección de dos alturas; la
      perpendicularidad se impone con el producto punto.
  \end{enumerate}

  \vspace{0.3cm}
  \destaca{\small \textit{Resultados para verificar al final.}
  Circuncentro $\left(3,\tfrac{7}{3}\right)$; baricentro
  $\left(\tfrac{8}{3},2\right)$; ortocentro
  $\left(2,\tfrac{4}{3}\right)$. La mediatriz de $AB$ es $x=3$; la de
  $AC$ es $x+3y=10$; la altura desde $C$ es $x=2$; la altura desde $A$
  es $2x-3y=0$.}
\end{frame}
% ============================================================
\begin{frame}{La recta de Euler}
  Con los tres puntos calculados, comprobemos algo notable:
  \[
    \underbrace{(3,\tfrac73)}_{\text{circuncentro}}\qquad
    \underbrace{(\tfrac83,2)}_{\text{baricentro}}\qquad
    \underbrace{(2,\tfrac43)}_{\text{ortocentro}}
  \]
  El vector del circuncentro al ortocentro es $(-1,-1)$; el vector del
  circuncentro al baricentro es $\left(-\tfrac13,-\tfrac13\right)$.
  Uno es múltiplo del otro, así que los tres puntos son
  \textbf{colineales}, y el baricentro está a un tercio del camino.

  \vspace{0.3cm}
  \destaca{Esta recta se llama \textbf{recta de Euler}, y el
  alineamiento ocurre en \textbf{todo} triángulo (no equilátero). Lo
  verificamos hoy con las herramientas de estas tres semanas
  (coordenadas, vectores y producto punto). El reto mayor, para quien
  lo quiera: demostrarlo en general, con un triángulo $A$, $B$, $C$
  arbitrario.}
\end{frame}
% ============================================================
\begin{frame}{Cierre de la semana}
  \textbf{Lo que vimos}
  {\small
  \begin{itemize}
    \item Ángulo entre rectas por dos métodos; proyección de un vector
      sobre otro
    \item La distancia punto--recta, y la forma normal clásica como la
      misma fórmula en otra notación
    \item Familias de rectas con un parámetro: el parámetro como
      selector
    \item Los puntos notables del triángulo y la recta de Euler
  \end{itemize}}

  \vspace{0.2cm}
  \textbf{Para la próxima semana}
  {\small
  \begin{itemize}
    \item Cuestionario de práctica de la semana 3 en Moodle
      (intentos ilimitados; no califica)
    \item \textbf{Tarea de redacción} en Moodle (la fecha límite
      aparece en la plataforma)
    \item Lehmann: series sobre la forma normal, la distancia
      punto--recta y las familias de rectas
    \item La próxima semana hay laboratorio: \textbf{Construcción 2}
  \end{itemize}}

  \vspace{0.2cm}
  \destaca{\textbf{Pregunta para la próxima sesión:} $\pp{u}{v}=0$
  indica perpendicularidad. ¿Qué información proporciona
  $\pp{u}{v}<0$? Y el anuncio de la semana 4: cierre de la Parte I
  (repaso y Examen 1) y la Construcción 2 en el laboratorio.}
\end{frame}
% ============================================================
%  Página de cierre del archivo.
%  Ajustar la URL cuando el libro quede publicado en GitHub Pages.
\begin{frame}[plain]
  \centering

  \vspace{1.0cm}
  \begin{tikzpicture}[scale=0.6,>=Stealth]
    \draw[very thin,color=gray!30] (-0.9,-0.9) grid (4.1,2.9);
    \draw[->,thick] (-0.9,0) -- (4.3,0);
    \draw[->,thick] (0,-0.9) -- (0,3.1);
    \draw[dorado,very thick,->] (0,0) -- (3.6,0.9);
    \draw[terracota,very thick,->] (0,0) -- (1.4,2.3);
    \coordinate (S) at ({(1.4*3.6+2.3*0.9)/(3.6*3.6+0.9*0.9)*3.6},
                        {(1.4*3.6+2.3*0.9)/(3.6*3.6+0.9*0.9)*0.9});
    \draw[dashed,gray,thick] (1.4,2.3) -- (S);
    \draw[turquesaGA,ultra thick] (0,0) -- (S);
  \end{tikzpicture}

  \vspace{0.5cm}

  {\Large\color{turquesaGA} ¿Preguntas?}\\[10pt]
  {\small\color{grisT}
  \url{https://rosaliahdez.github.io/notas-geometria-analitica/} \\
  Comunicación por Teams · Cuestionarios y entregas en Moodle (Ramanujan)}

  \vspace{0.7cm}

  {\tiny\color{grisT}
  Texto base: Ramírez Galarza, \textit{Geometría analítica. Una
  introducción a la geometría}, UNAM. ---
  Ejercicios: Lehmann, \textit{Geometría analítica}, Limusa.}
\end{frame}
\end{document}
